
\chapter{Modelli lineari: i metodi di soluzione iterativi}
\newpage
\minitoc
\newpage

Il modello presentato in questo documento è lo stesso trattato nel capitolo 3 del libro ``Computaional Economics. A concise introduction'' di Afonso e Vasconcelos pubblicato nel 2016 dalla casa editrice Routledge.  

\section{Il modello IS--LM in economia aperta e i metodi iterativi}
\label{sec:islm-aperta-iterativi}

Nella parte relativa ai metodi di soluzione diretti abbiamo utilizzato il modello
IS--LM in economia chiusa per introdurre la fattorizzazione LU.
La fattorizzazione LU appartiene alla famiglia dei
\emph{metodi diretti}: una volta assegnati la matrice $A$ e il
vettore $b$, la soluzione del sistema

\begin{equation*}
    Ax=b
\end{equation*}

viene ottenuta attraverso una sequenza finita e predeterminata di
operazioni.

Consideriamo ora un'estensione del modello IS--LM nella quale due
economie sono collegate attraverso gli scambi internazionali.
L'aumento del numero delle equazioni e delle variabili ci offre
l'occasione per introdurre una diversa famiglia di tecniche
numeriche: i \emph{metodi iterativi}.

A differenza dei metodi diretti, i metodi iterativi partono da una
soluzione iniziale di prova e costruiscono una successione di
approssimazioni della soluzione. Questi metodi risultano
particolarmente interessanti quando si devono affrontare sistemi
caratterizzati da matrici grandi e sparse, cioè matrici nelle quali
molti elementi sono nulli.


\subsection{Il modello IS--LM in economia aperta}

Consideriamo due paesi. Le variabili relative al secondo paese
sono indicate con il pedice $f$.

Per il primo paese abbiamo

\begin{align*}
    Y &= C+I+G+NX,\\
    C &= \bar C+c(Y-T),\\
    I &= \bar I-bR,\\
    NX &= \overline{NX}-j(Y-Y_f)
          +lE\frac{\bar P_f}{P},\\
    \frac{M}{P} &= kY-hR.
\end{align*}

Per il secondo paese abbiamo invece

\begin{align*}
    Y_f &= C_f+I_f+G_f+NX_f,\\
    C_f &= \bar C_f+c_f(Y_f-T_f),\\
    I_f &= \bar I_f-b_fR_f,\\
    NX_f &= -NX,\\
    \frac{M_f}{P_f} &= k_fY_f-h_fR_f.
\end{align*}

Rispetto al modello in economia chiusa compare ora un nuovo
elemento: le esportazioni nette $NX$ collegano le due economie.

La relazione

\begin{equation*}
    NX_f=-NX
\end{equation*}

esprime il fatto che le esportazioni nette di un paese
corrispondono alle importazioni nette dell'altro.

Abbiamo infine l'ulteriore condizione di parità dei tassi di interesse:
\[R=R_f\]

Abbiamo dunque in totale 11 equazioni. Le variabili endogene sono allo stesso modo 11: $Y,\ C,\ I,\ NX,\ R,\ Y_f,\ C_f,\ I_f,\ NX_f,\ R_f,\ E$.

I parametri e le variabili esogene sono riportati nella seguente tabella:

\begin{center}
\scriptsize
\renewcommand{\arraystretch}{1.15}
\begin{tabular}{ll|ll}
\hline
\multicolumn{2}{c|}{\textbf{Economia nazionale}} &
\multicolumn{2}{c}{\textbf{Economia estera}}\\
\hline
\multicolumn{4}{c}{\textbf{Parametri}}\\
\hline
$c=0.63$ & Prop. al consumo &
$c_f=0.63$ & Prop. al consumo\\
$b=1500$ & Sens. investimenti &
$b_f=1500$ & Sens. investimenti\\
$k=0.6$ & Sens. moneta al reddito &
$k_f=0.6$ & Sens. moneta al reddito\\
$h=2700$ & Sens. moneta al tasso &
$h_f=2700$ & Sens. moneta al tasso\\
\hline
\multicolumn{4}{c}{\textbf{Parametri degli scambi}}\\
\hline
\multicolumn{2}{l|}{$j=0.1$: sensibilità ai redditi} &
\multicolumn{2}{l}{$l=10$: sensibilità al cambio reale}\\
\hline
\multicolumn{4}{c}{\textbf{Variabili esogene}}\\
\hline
$\bar C=55$ & Consumo autonomo &
$\bar C_f=55$ & Consumo autonomo\\
$\bar I=47$ & Investimento autonomo &
$\bar I_f=47$ & Investimento autonomo\\
$\overline{NX}=-10$ & Export netto autonomo &
&\\
$\bar G=150$ & Spesa pubblica &
$\bar G_f=150$ & Spesa pubblica\\
$\bar M=210$ & Offerta di moneta &
$\bar M_f=210$ & Offerta di moneta\\
$\bar T=150$ & Imposte &
$\bar T_f=150$ & Imposte\\
$\bar P=1$ & Livello dei prezzi &
$\bar P_f=1$ & Livello dei prezzi\\
\hline
\end{tabular}
\end{center}

Date le variabili esogene e i parametri, il problema consiste
ancora una volta nel determinare simultaneamente le variabili
endogene del modello. Una volta ottenuto l'equilibrio, possiamo
inoltre modificare gli strumenti di politica fiscale o monetaria
e studiarne gli effetti sulle due economie.

Anche questo modello può essere ricondotto alla forma

\begin{equation*}
    Ax=b.
\end{equation*}
dove:

\setcounter{MaxMatrixCols}{11}
\tiny
$$
A=\begin{pmatrix}
1&-1&-1&-1&0&0&0&0&0&0&0\\
-c&1&0&0&0&0&0&0&0&0&0\\
0&0&1&0&b&0&0&0&0&0&0\\
k&0&0&0&-h&0&0&0&0&0&0\\
0&0&0&0&0&1&-1&-1&-1&0&0\\
0&0&0&0&0&-c_f&1&0&0&0&0\\
0&0&0&0&0&0&0&1&0&b_f&0\\
0&0&0&0&0&k_f&0&0&0&-h_f&0\\
j&0&0&1&0&-j&0&0&0&0&-l\frac{\bar P_f}{\bar P}\\
0&0&0&1&0&0&0&0&1&0&0\\
0&0&0&0&1&0&0&0&0&-1&0
\end{pmatrix}
x=\begin{pmatrix}
Y\\ C\\ I\\ NX\\ R\\ Y_f\\ C_f\\ I_f\\ NX_f\\ R_f\\ E
\end{pmatrix}
b=\begin{pmatrix}
\bar G\\ \bar C-c\bar T\\ \bar I\\ \frac{\bar M}{\bar P}\\ \bar G_f\\ \bar C_f-c_f\bar T_f\\ \bar I_f\\ \frac{\bar M_f}{\bar P_f}\\ \overline{NX}\\ 0\\ 0
\end{pmatrix}
$$

\normalsize

Il problema matematico rimane quindi quello già incontrato nelle
sezioni precedenti. Cambia, tuttavia, la tecnica numerica che
vogliamo utilizzare per determinarne la soluzione.


\subsection{L'idea dei metodi iterativi}

Un metodo iterativo non cerca necessariamente di ottenere la
soluzione attraverso una sequenza prefissata di operazioni.
Partendo da un vettore iniziale $x_0$, genera invece una successione

\begin{equation*}
    x_0,\;x_1,\;x_2,\;\ldots
\end{equation*}

che, se il metodo converge, si avvicina progressivamente alla
soluzione del sistema.

Per costruire un meccanismo iterativo partiamo ancora da

\begin{equation*}
    Ax=b.
    \label{eq:iter-system}
\end{equation*}

Introduciamo una matrice $M$ e scriviamo

\begin{equation*}
    A=M-(M-A).
\end{equation*}

Sostituendo questa espressione nel sistema otteniamo

\begin{equation*}
    \left[M-(M-A)\right]x=b,
\end{equation*}

da cui

\begin{equation*}
    Mx=(M-A)x+b.
\end{equation*}

Introducendo gli indici di iterazione, possiamo utilizzare i valori
calcolati all'iterazione $k$ per costruire una nuova approssimazione:

\begin{equation*}
    Mx_{k+1}=(M-A)x_k+b.
\end{equation*}

Sottraendo $Mx_k$ da entrambi i membri otteniamo

\begin{equation}
    M(x_{k+1}-x_k)=b-Ax_k.
    \label{eq:iter-general}
\end{equation}

Questa relazione costituisce la base del meccanismo iterativo.


\subsection{Correzione e residuo}

Per interpretare l'Equazione~\eqref{eq:iter-general}, definiamo

\begin{equation*}
    z_k=x_{k+1}-x_k
    \label{eq:correction}
\end{equation*}

e

\begin{equation*}
    r_k=b-Ax_k.
    \label{eq:residual}
\end{equation*}

Il vettore $z_k$ rappresenta la \emph{correzione} che deve essere
aggiunta all'approssimazione corrente per ottenere quella
successiva.

Il vettore $r_k$ prende invece il nome di \emph{residuo}. Esso
misura quanto l'approssimazione corrente $x_k$ non soddisfa ancora
il sistema originale.

Se $x_k$ fosse già la soluzione esatta, avremmo infatti

\begin{equation*}
    Ax_k=b
\end{equation*}

e quindi

\begin{equation*}
    r_k=0.
\end{equation*}

Utilizzando queste definizioni, l'Equazione~\eqref{eq:iter-general}
diventa

\begin{equation*}
    Mz_k=r_k.
    \label{eq:Mzr}
\end{equation*}

Il problema iterativo può quindi essere interpretato in modo molto
semplice: calcoliamo l'errore commesso dall'approssimazione corrente,
espresso dal residuo, e utilizziamo questa informazione per
determinare una correzione $z_k$.


\subsection{Il meccanismo iterativo generale}

Possiamo ora descrivere la procedura completa.

Si parte scegliendo arbitrariamente un vettore iniziale

\begin{equation*}
    x_0.
\end{equation*}

Si calcola il residuo

\begin{equation*}
    r_0=b-Ax_0.
\end{equation*}

Successivamente si risolve il sistema

\begin{equation*}
    Mz_0=r_0
\end{equation*}

per determinare la correzione $z_0$.

La nuova approssimazione è quindi

\begin{equation*}
    x_1=x_0+z_0.
\end{equation*}

La stessa procedura viene ripetuta:

\begin{equation*}
    r_1=b-Ax_1,
\end{equation*}

\begin{equation*}
    Mz_1=r_1,
\end{equation*}

\begin{equation*}
    x_2=x_1+z_1,
\end{equation*}

e così via.

In termini generali, a ogni iterazione $k$ eseguiamo quindi i tre
passaggi

\begin{align*}
    r_k &= b-Ax_k,\\
    Mz_k &= r_k,\\
    x_{k+1} &= x_k+z_k.
\end{align*}

Questo schema costituisce il meccanismo iterativo generale.
%denominato \emph{metodo di Richardson}.


\subsection{La scelta della matrice $M$}

A ogni iterazione dobbiamo risolvere il sistema

\begin{equation*}
    Mz_k=r_k.
\end{equation*}

Questa operazione rappresenta la parte computazionalmente più
impegnativa del procedimento.

Se la soluzione di questo sistema fosse difficile quanto quella
del sistema originale $Ax=b$, il metodo iterativo non offrirebbe
un particolare vantaggio. La matrice $M$ deve quindi essere scelta
in modo tale che il sistema $Mz_k=r_k$ sia molto semplice da
risolvere.

Due scelte particolarmente convenienti sono una matrice diagonale
oppure una matrice triangolare.

Scelte differenti di $M$ danno origine a differenti metodi
iterativi.

Nel \emph{metodo di Jacobi} si utilizza

\begin{equation*}
    M=\operatorname{diag}(A),
\end{equation*}

cioè la matrice contenente gli elementi della diagonale di $A$ e
zeri in tutte le altre posizioni.

Nel \emph{metodo di Gauss--Seidel}, invece, $M$ è costituita dalla
parte triangolare inferiore di $A$.

Poiché una matrice diagonale o triangolare rende particolarmente
semplice la soluzione di $Mz_k=r_k$, ciascuna iterazione può essere
eseguita con un costo relativamente contenuto.

\`E possibile introdurre inoltre il \emph{metodo del rilassamento},
nel quale la diagonale di $M$ viene modificata attraverso un
parametro $\omega\neq0$.


\subsection{Convergenza dei metodi stazionari}

Il metodo iterativo non garantisce che la
successione

\begin{equation*}
    x_0,x_1,x_2,\ldots
\end{equation*}

converga effettivamente alla soluzione.

Le condizioni di convergenza dipendono dal metodo e dalle
caratteristiche della matrice $A$.

Per il metodo di Jacobi, la condizione è
che $A$ sia a dominanza diagonale stretta per righe.

Per Gauss--Seidel, ci sono due possibilità:

\begin{itemize}
    \item $A$ a dominanza diagonale stretta per righe;
    \item $A$ simmetrica e definita positiva.
\end{itemize}

Queste restrizioni rappresentano un limite importante dei metodi
iterativi considerati fino a questo punto.


\subsection{Introduzione di un parametro di passo}

Il meccanismo precedente aggiorna la soluzione secondo la regola

\begin{equation*}
    x_{k+1}=x_k+z_k.
\end{equation*}

Una possibile estensione consiste nell'introdurre un parametro
$\alpha_k$ nell'equazione precedente trasformandola nella seguente:

\begin{equation*}
    x_{k+1}=x_k+\alpha_k z_k
    \label{eq:iter-alpha}
\end{equation*}

Il termine $\alpha_kz_k$ determina quindi il movimento effettuato
a ogni iterazione.

Possiamo interpretare separatamente le due componenti:

\begin{itemize}
    \item $z_k$ determina la \emph{direzione} del movimento;
    \item $\alpha_k$ determina l'\emph{intensità} del movimento.
\end{itemize}

Per ottenere una convergenza rapida è quindi importante scegliere
in modo appropriato sia la direzione sia l'ampiezza del passo.


\subsection{Metodi stazionari e non stazionari}

La presenza del parametro $\alpha_k$ permette di introdurre una
distinzione tra due famiglie di metodi iterativi.

Nei \emph{metodi stazionari}, il parametro non cambia nel corso
delle iterazioni:

\begin{equation*}
    \alpha_k=\alpha.
\end{equation*}

I metodi di Jacobi e Gauss--Seidel appartengono alla famiglia dei
metodi stazionari.

Nei \emph{metodi non stazionari}, invece, il parametro può cambiare
a ogni iterazione: $\alpha_k$ 
viene quindi scelto in funzione dell'informazione disponibile al
passo $k$.

I metodi stazionari hanno il vantaggio di essere semplici, ma
convergono soltanto sotto determinate condizioni. Questo motiva
l'introduzione di metodi nei quali direzione e intensità del
movimento possono essere adattate nel corso delle iterazioni.


\subsection{Il caso di una matrice simmetrica e definita positiva}

Consideriamo ora il caso in cui $A$ sia simmetrica e definita
positiva.

Associamo al sistema $Ax=b$ la funzione

\begin{equation*}
    \Phi(x)
    =
    \frac{1}{2}x^TAx-x^Tb.
    \label{eq:phi-gradient}
\end{equation*}

Il gradiente di questa funzione è

\begin{equation}
    \nabla\Phi(x)=Ax-b.
    \label{eq:gradient-phi}
\end{equation}

Se $A$ è simmetrica e definita positiva, $\Phi(x)$ è strettamente
convessa e possiede quindi un unico punto di minimo.

Nel punto di minimo il gradiente è nullo:

\begin{equation*}
    \nabla\Phi(x)=0.
\end{equation*}

Utilizzando l'Equazione~\eqref{eq:gradient-phi}, questa condizione
diventa

\begin{equation*}
    Ax-b=0,
\end{equation*}

ovvero

\begin{equation*}
    Ax=b.
\end{equation*}

La soluzione di un sistema lineare con matrice simmetrica e
definita positiva può quindi essere interpretata anche come il
punto che minimizza la funzione quadratica $\Phi(x)$.

Abbiamo così trasformato il problema di soluzione del sistema in
un problema di minimizzazione.


\subsection{Il gradiente come direzione}

Per costruire un metodo iterativo dobbiamo ora scegliere una
direzione $z_k$ che ci permetta di avvicinarci al minimo della
funzione $\Phi$.

Il gradiente

\begin{equation*}
    \nabla\Phi(x)
\end{equation*}

indica la direzione di massima crescita della funzione. Per
comprendere questa proprietà consideriamo la derivata direzionale
di una generica funzione $f$ nel punto $x_0$, lungo una direzione
rappresentata da un vettore avente norma euclidea uguale a 1 (possiamo parlare in questo caso di versore) denotato con $u$:

\begin{equation*}
    D_uf(x_0)=\nabla f(x_0)\cdot u,
    \qquad
    \|u\|=1.
\end{equation*}

Utilizzando la formula geometrica del prodotto scalare possiamo scrivere

\begin{equation*}
    \nabla f(x_0)\cdot u
    =
    \|\nabla f(x_0)\|\,\|u\|\cos\theta,
\end{equation*}

dove $\theta$ è l'angolo tra $\nabla f(x_0)$ e $u$.

Poiché $u$ è un versore,

\begin{equation*}
    \|u\|=1,
\end{equation*}

e quindi

\begin{equation*}
    D_uf(x_0)
    =
    \|\nabla f(x_0)\|\cos\theta.
\end{equation*}

Il valore massimo si ottiene quando

\begin{equation*}
    \cos\theta=1,
\end{equation*}

cioè quando $\theta=0$. In questo caso $u$ ha la stessa direzione
del gradiente:

\begin{equation*}
    u=
    \frac{\nabla f(x_0)}
         {\|\nabla f(x_0)\|}.
\end{equation*}

La massima derivata direzionale è pertanto

\begin{equation*}
    \max_{\|u\|=1}D_uf(x_0)
    =
    \|\nabla f(x_0)\|.
\end{equation*}

Il gradiente indica dunque la direzione di massima crescita.
Se vogliamo invece ridurre $\Phi$, dobbiamo muoverci nel verso
opposto.

Una possibile scelta della direzione è quindi

\begin{equation*}
    z_k=-\nabla\Phi(x_k).
    \label{eq:steepest-direction}
\end{equation*}


\subsection{La scelta dell'intensità del movimento}

Una volta individuata la direzione $z_k$, rimane da stabilire
quanto muoversi lungo quella direzione.

Partendo dalla regola

\begin{equation*}
    x_{k+1}=x_k+\alpha_kz_k,
\end{equation*}

possiamo scegliere $\alpha_k$ in modo da minimizzare la funzione
lungo la direzione $z_k$:

\begin{equation*}
    \alpha_k
    =
    \arg\min_{\alpha}
    \Phi(x_k+\alpha z_k).
    \label{eq:line-search}
\end{equation*}

La combinazione

\begin{equation*}
    z_k=-\nabla\Phi(x_k)
\end{equation*}

con la scelta dell'ampiezza del passo conduce al
\emph{metodo del gradiente}.

Quando $A$ è simmetrica e definita positiva, il metodo converge
alla soluzione a partire da qualsiasi punto iniziale $x_0$.

Un'evoluzione di questa idea è il \emph{metodo del gradiente
coniugato}, nel quale vengono costruite direzioni che permettono
di accelerare la convergenza.


\subsection{Il caso di una matrice generica}

La situazione è più complessa quando $A$ non è simmetrica e
definita positiva.

La soluzione iterativa di sistemi con matrici generiche costituisce
un argomento molto ampio. In questo corso non ne svilupperemo in
dettaglio la teoria, ma utilizzeremo uno dei metodi maggiormente
diffusi: il \emph{Generalized Minimal Residual Method}, indicato
con l'acronimo GMRES.

Il metodo GMRES appartiene alla famiglia dei metodi che costruiscono
le nuove approssimazioni cercando di ridurre il residuo.

Nelle applicazioni utilizzeremo direttamente le implementazioni
disponibili nei principali software scientifici. Ad esempio:

\begin{verbatim}
# R
pracma::gmres(A, b, ...)

# Octave
gmres(A, b, ...)

# Python
scipy.sparse.linalg.gmres(A, b, ...)

# Julia
IterativeSolvers.gmres(A, b, ...)
\end{verbatim}

L'obiettivo, in questa fase, non è ricostruire l'algoritmo GMRES,
ma comprendere perché possa essere necessario ricorrere a un metodo
iterativo più generale quando le proprietà della matrice non
permettono di utilizzare convenientemente i metodi precedenti.


\subsection{Una classificazione dei metodi per sistemi lineari}

Possiamo a questo punto riordinare le tecniche introdotte finora
in due grandi famiglie.

I \emph{metodi diretti} determinano la soluzione attraverso un
numero predefinito di operazioni. In questa categoria abbiamo
incontrato:

\begin{itemize}
    \item l'eliminazione di Gauss;
    \item la fattorizzazione LU.
\end{itemize}

I \emph{metodi iterativi}, invece, costruiscono una successione
di approssimazioni della soluzione. Possiamo distinguerli
ulteriormente in metodi stazionari e non stazionari.

Tra i metodi stazionari abbiamo considerato:

\begin{itemize}
    \item lo schema generale di Richardson;
    \item il metodo di Jacobi;
    \item il metodo di Gauss--Seidel;
    \item il metodo del rilassamento.
\end{itemize}

Tra i metodi non stazionari, la scelta dipende anche dalle
caratteristiche della matrice. Per matrici simmetriche e definite
positive abbiamo introdotto:

\begin{itemize}
    \item il metodo del gradiente;
    \item il metodo del gradiente coniugato.
\end{itemize}

Per matrici generiche abbiamo invece introdotto il metodo

\begin{equation*}
    \text{GMRES}.
\end{equation*}

La scelta del metodo numerico non è quindi indipendente dalla
struttura del problema. Le dimensioni della matrice, la presenza
di molti elementi nulli e proprietà quali simmetria, definita
positività e dominanza diagonale possono rendere un metodo più
appropriato di un altro.


\subsection{Il ruolo del residuo nei metodi iterativi}

Un elemento comune ai metodi iterativi è il residuo

\begin{equation*}
    r_k=b-Ax_k.
\end{equation*}

Esso fornisce informazioni sulla qualità dell'approssimazione
corrente e può essere utilizzato sia per determinare come
procedere verso una nuova approssimazione, sia per stabilire
quando interrompere il procedimento.

I metodi iterativi più avanzati costruiscono le nuove
approssimazioni all'interno di particolari sottospazi, detti
\emph{sottospazi di Krylov}, che si ampliano progressivamente
con il procedere delle iterazioni.

All'interno di questa famiglia possono essere distinti due
approcci.

Nel primo, detto \emph{Ritz--Galerkin} o del residuo ortogonale,
la nuova approssimazione viene scelta imponendo una condizione
di ortogonalità sul residuo. A questa famiglia appartengono il
gradiente coniugato, utilizzato per matrici simmetriche e definite
positive, e il bi-CG per matrici generiche.

Nel secondo, detto \emph{approccio del minimo residuo}, si sceglie
invece l'approssimazione che rende minima l'ampiezza del residuo
all'interno del sottospazio considerato. Il metodo GMRES appartiene
a questa seconda famiglia.

Questi concetti forniscono un'indicazione della logica sulla quale
si basano i metodi iterativi moderni. In questo corso non
svilupperemo ulteriormente la teoria dei sottospazi di Krylov, ma
utilizzeremo gli algoritmi disponibili nei software di calcolo
scientifico quando le applicazioni lo richiederanno.


\subsection{Dal problema economico alla soluzione iterativa}

Possiamo infine sintetizzare il percorso seguito in questa sezione:

\begin{equation*}
\boxed{
\begin{array}{c}
\text{modello IS--LM in economia aperta}\\
\downarrow\\
\text{sistema lineare }Ax=b\\
\downarrow\\
\text{scelta di un'approssimazione iniziale }x_0\\
\downarrow\\
r_k=b-Ax_k\\
\downarrow\\
\text{determinazione della correzione}\\
\downarrow\\
x_{k+1}=x_k+\alpha_kz_k\\
\downarrow\\
\text{nuova approssimazione della soluzione}
\end{array}}
\end{equation*}

Il passaggio dai metodi diretti ai metodi iterativi introduce
quindi un modo diverso di affrontare il problema $Ax=b$.
Invece di costruire direttamente la soluzione, si genera una
sequenza di approssimazioni che viene progressivamente corretta
utilizzando l'informazione contenuta nel residuo.

\subsection{Soluzione del modello IS-LM con economia aperta}

Per determinare i valori di equilibrio delle variabili
endogene, abbiamo utilizzato la funzione
\texttt{gmres()} del pacchetto R \texttt{pracma}.
La funzione implementa il metodo GMRES
(\emph{Generalized Minimal Residual}).
A differenza dei metodi diretti, GMRES costruisce
una successione di approssimazioni della soluzione,
minimizzando a ogni iterazione la norma euclidea
del residuo in un opportuno sottospazio di Krylov.
Il residuo è definito come
\[
\mathbf{r}_k=\mathbf{d}-A\mathbf{x}_k,
\]
dove $\mathbf{x}_k$ rappresenta l'approssimazione
della soluzione alla $k$-esima iterazione.

La sintassi della funzione è:

\begin{verbatim}
gmres(A, b, x0 = rep(0, length(b)),
      errtol = 1e-6,
      kmax = length(b) + 1,
      reorth = 1)
\end{verbatim}

I primi due argomenti sono obbligatori.
L'argomento \texttt{A} rappresenta la matrice
quadrata dei coefficienti del sistema, mentre
\texttt{b} è il vettore dei termini noti, che nel
nostro modello abbiamo denominato \texttt{d}.

Gli altri argomenti sono facoltativi.
\texttt{x0} specifica il vettore iniziale da cui
parte il procedimento iterativo; per impostazione
predefinita è il vettore nullo.
\texttt{errtol} stabilisce la tolleranza relativa
per la riduzione del residuo, con valore
predefinito pari a $10^{-6}$.
\texttt{kmax} indica il numero massimo di
iterazioni, mentre \texttt{reorth} controlla
la procedura di riortogonalizzazione dei
vettori utilizzati dall'algoritmo.
Il valore predefinito di quest'ultimo parametro
è 1, corrispondente al criterio di
Brown--Hindmarsh.

Nel nostro caso, una volta costruiti la matrice
\texttt{A} e il vettore \texttt{d}, il sistema
può essere risolto con le istruzioni:

\begin{verbatim}
library(pracma)
risultato <- gmres(A, d)
\end{verbatim}

La funzione restituisce una lista contenente
tre elementi. L'elemento \texttt{x} è il vettore
della soluzione approssimata del sistema;
\texttt{error} contiene la successione delle
norme dei residui calcolate durante il procedimento;
\texttt{niter} indica il numero di iterazioni
effettuate.

Possiamo quindi estrarre e visualizzare
separatamente i risultati:

\begin{verbatim}
soluzione <- risultato$x
errori <- risultato$error
iterazioni <- risultato$niter
\end{verbatim}

Il vettore \texttt{soluzione} contiene i valori
di equilibrio delle undici variabili endogene,
nell'ordine stabilito nella costruzione della
matrice dei coefficienti.

Per verificare la precisione numerica della
soluzione ottenuta, possiamo inoltre calcolare
direttamente la norma euclidea del residuo:

\begin{verbatim}
residuo <- d - as.vector(A %*% soluzione)
norma_residuo <- sqrt(sum(residuo^2))
\end{verbatim}

Un valore sufficientemente piccolo della norma
del residuo indica che la soluzione approssimata
soddisfa il sistema lineare con buona precisione
numerica.

La soluzione di equilibrio del modello è riportata nella prima colonna della tabella \ref{tab:eqislm}.

Una volta determinata la soluzione di equilibrio,
è possibile analizzare gli effetti di interventi di politica
economica modificando le variabili esogene e calcolando il
nuovo equilibrio del sistema. Consideriamo due esperimenti
alternativi di politica fiscale espansiva, entrambi riferiti
all'economia nazionale: nel primo si ipotizza una riduzione
delle imposte di 20 unità ($\Delta \bar{T}=-20$), mentre
nel secondo si considera un aumento della spesa pubblica
di 20 unità ($\Delta \bar{G}=+20$).
In ciascun esperimento, tutte le altre variabili esogene
rimangono invariate.

La tabella riporta, per ciascuna delle undici variabili
endogene, il valore di equilibrio iniziale nella colonna
\emph{Soluzione} e le variazioni rispetto a tale equilibrio
prodotte dai due interventi nelle colonne $\Delta T$ e
$\Delta G$. Queste ultime non rappresentano quindi i nuovi
livelli di equilibrio, ma le differenze tra i valori
ottenuti dopo l'intervento e quelli della soluzione iniziale.

Entrambe le politiche determinano un aumento del reddito
nazionale ed estero, accompagnato da un incremento dei
tassi di interesse e da una riduzione degli investimenti.
L'aumento della spesa pubblica produce un incremento del
reddito nazionale maggiore rispetto alla riduzione delle
imposte ($14.440$ contro $8.957$), mentre la riduzione
delle imposte determina un aumento più consistente del
consumo nazionale ($18.243$ contro $9.097$).
In entrambi i casi, le esportazioni nette nazionali
diminuiscono e quelle estere aumentano, mentre il tasso
di cambio $E$ si riduce.

\begin{table}
\small
\renewcommand{\arraystretch}{1.15}
\setlength{\tabcolsep}{9pt}
\centering
\begin{tabular}{lrrr}
\hline
\textbf{Variabile} &
\textbf{Soluzione} &
\textbf{$\Delta T$} &
\textbf{$\Delta G$} \\
\hline
$Y$    & 389.810 &  8.957 & 14.440 \\
$C$    & 206.081 & 18.243 &  9.097 \\
$I$    &  33.730 & -2.986 & -4.813 \\
$NX$   &   0.000 & -6.300 & -9.844 \\
$R$    & 0.008847 & 0.001991 & 0.003209 \\
$Y_f$  & 389.810 &  8.957 & 13.996 \\
$C_f$  & 206.081 &  5.643 &  8.817 \\
$I_f$  &  33.730 & -2.986 & -4.665 \\
$NX_f$ &   0.000 &  6.300 &  9.844 \\
$R_f$  & 0.008847 & 0.001991 & 0.003110 \\
$E$    &   1.000 & -0.630 & -0.980 \\
\hline
\end{tabular}
\caption{Valori di equilibrio e variazioni di tali valori a seguito di riduzioni dell'imposizione di 20 e di aumento della spesa pubblica di 20}
\label{tab:eqislm}
\end{table}


\section{Implementazione dei metodi iterativi}
\label{sec:implementazione-metodi-iterativi}


Nella sezione precedente abbiamo utilizzato la funzione
\texttt{gmres()} di R, che esegue automaticamente le operazioni necessarie per ottenere la soluzione. 
Tuttavia, utilizzare una funzione già disponibile
è diverso dal comprendere come programmare un procedimento
iterativo. Quest'ultima competenza è particolarmente
importante quando si affrontano modelli economici per i quali
non esistono funzioni predefinite adatte alla loro soluzione,
oppure quando è necessario costruire algoritmi numerici
specifici.

Ricordiamo che i metodi iterativi sono utili a trovare 
la soluzione di un sistema lineare
\begin{equation*}
    Ax=b.
\end{equation*}
quando la matrice $A$ e grande e sparsa.

L'idea fondamentale consiste nel partire da un'approssimazione
iniziale $x_0$ e costruire una successione

\begin{equation*}
    x_0,\;x_1,\;x_2,\;\ldots
\end{equation*}

che, sotto opportune condizioni, converge alla soluzione del sistema.

In termini generali, a ogni iterazione viene calcolato il residuo

\begin{equation*}
    r_k=b-Ax_k,
\end{equation*}

dal quale si ricava una correzione della soluzione corrente.
Il procedimento viene quindi ripetuto fino a quando
l'approssimazione ottenuta può essere considerata sufficientemente
accurata.

Dal punto di vista teorico, il meccanismo è relativamente semplice.
Tuttavia, qualora dovessimo implementare l'algoritmo dall'inizio senza ricorrere a delle funzioni già disponibili avremmo due nuovi problemi computazionali:

\begin{enumerate}
    \item come possiamo chiedere al computer di ripetere più volte
          la stessa sequenza di operazioni?
    \item come possiamo stabilire quando il procedimento iterativo
          deve essere interrotto?
\end{enumerate}

Per rispondere alla prima domanda introduciamo i cicli
\texttt{for} e \texttt{while}. Per rispondere alla seconda
introduciamo il concetto di norma e lo utilizziamo per misurare
la dimensione del residuo.


\subsection{Ripetere un'operazione: il ciclo \texttt{for}}

Un algoritmo iterativo richiede di eseguire ripetutamente la stessa
sequenza di operazioni.

Supponiamo, ad esempio, di voler eseguire un'operazione per un numero
prefissato di volte. In questo caso possiamo utilizzare un ciclo
\texttt{for}.

La struttura generale in \textsf{R} è

\begin{verbatim}
for (i in sequenza) {

    operazioni

}
\end{verbatim}

La variabile $i$ prende, uno dopo l'altro, i valori contenuti nella
sequenza. Per ciascun valore vengono eseguite le istruzioni
contenute nel ciclo.

Ad esempio,

\begin{verbatim}
for (i in 1:10) {

    print(i)

}
\end{verbatim}

esegue il comando \texttt{print(i)} dieci volte, assegnando
progressivamente a \texttt{i} i valori da 1 a 10.

Nel contesto di un metodo iterativo, l'indice $i$ può essere
interpretato come l'indice dell'iterazione:

\begin{equation*}
    i=1,2,\ldots,K,
\end{equation*}

dove $K$ rappresenta il numero massimo di iterazioni che vogliamo
eseguire.

Una struttura di questo tipo permette quindi di tradurre in codice
una successione del tipo

\begin{equation*}
    x_0
    \longrightarrow
    x_1
    \longrightarrow
    x_2
    \longrightarrow
    \cdots
    \longrightarrow
    x_K.
\end{equation*}


\subsection{Ripetere un'operazione finché una condizione è soddisfatta}

Nei metodi iterativi, tuttavia, non sempre sappiamo in anticipo
quante iterazioni saranno necessarie.

Potremmo voler continuare il procedimento fino a quando la soluzione
raggiunge un determinato livello di accuratezza.

In questo caso è più naturale utilizzare un ciclo \texttt{while}.

La struttura generale è

\begin{verbatim}
while (condizione) {

    operazioni

}
\end{verbatim}

Le istruzioni contenute nel ciclo vengono eseguite finché la
condizione specificata rimane vera.

Consideriamo, ad esempio,

\begin{verbatim}
i <- 1

while (i < 10) {

    print(i)
    i <- i + 1

}
\end{verbatim}

All'inizio abbiamo $i=1$. Dopo ogni iterazione il valore di $i$
viene aumentato di una unità. Il ciclo continua finché

\begin{equation*}
    i<10.
\end{equation*}

Quando questa condizione non è più verificata, il ciclo termina.

Il ciclo \texttt{while} è particolarmente adatto ai metodi
iterativi perché permette di formulare una regola del tipo:

\begin{quote}
continua a costruire nuove approssimazioni finché l'errore è
troppo grande.
\end{quote}

Per rendere operativa questa idea dobbiamo però definire come
misurare l'errore.


\subsection{Distanza e norma}

Per misurare quanto due numeri reali $a$ e $b$ siano diversi
possiamo utilizzare la distanza

\begin{equation*}
    d(a,b)=|a-b|.
\end{equation*}

Quando lavoriamo con vettori, la stessa idea richiede una
generalizzazione. Questa generalizzazione è fornita dal concetto
di \emph{norma}.

Una norma associa a un vettore un numero reale non negativo che
può essere interpretato come una misura della sua grandezza.

Per un vettore

\begin{equation*}
    x=
    \begin{pmatrix}
        x_1\\
        x_2\\
        \vdots\\
        x_n
    \end{pmatrix},
\end{equation*}

possiamo definire diverse norme.


\subsubsection{Norma 1}

La norma 1 è definita come

\begin{equation*}
    \|x\|_1
    =
    \sum_{i=1}^{n}|x_i|.
    \label{eq:norma1}
\end{equation*}

Essa corrisponde quindi alla somma dei valori assoluti degli
elementi del vettore.


\subsubsection{Norma 2}

La norma 2, detta anche norma euclidea, è

\begin{equation*}
    \|x\|_2
    =
    \sqrt{
        \sum_{i=1}^{n}x_i^2
    }.
    \label{eq:norma2}
\end{equation*}

Nel caso di un vettore con due componenti,

\begin{equation*}
    x=
    \begin{pmatrix}
        x_1\\
        x_2
    \end{pmatrix},
\end{equation*}

otteniamo

\begin{equation*}
    \|x\|_2
    =
    \sqrt{x_1^2+x_2^2}.
\end{equation*}

Questa espressione coincide con la distanza euclidea del punto
$(x_1,x_2)$ dall'origine.


\subsubsection{Norma infinito}

La norma infinito è invece definita come

\begin{equation*}
    \|x\|_\infty
    =
    \max_i |x_i|.
    \label{eq:normainf}
\end{equation*}

Essa misura quindi la componente del vettore che presenta il
valore assoluto maggiore.

Le tre norme forniscono misure differenti della grandezza di un
vettore. Nel seguito utilizzeremo principalmente la norma 2.


\subsection{La norma in \textsf{R}}

In \textsf{R}, la funzione \texttt{norm()} opera su oggetti di
tipo matrice. Se $x$ è inizialmente un vettore, possiamo quindi
trasformarlo in una matrice mediante

\begin{verbatim}
x <- as.matrix(x)
\end{verbatim}

e successivamente calcolarne la norma.

La norma 1 può essere ottenuta con

\begin{verbatim}
norm(x, type = "1")
\end{verbatim}

la norma 2 con

\begin{verbatim}
norm(x, type = "2")
\end{verbatim}

e la norma infinito con

\begin{verbatim}
norm(x, type = "I")
\end{verbatim}

La possibilità di calcolare numericamente la norma di un vettore
ci permette ora di costruire un criterio per decidere quando
interrompere un algoritmo iterativo.


\subsection{Il residuo come misura della qualità dell'approssimazione}

Riprendiamo il sistema

\begin{equation*}
    Ax=b.
\end{equation*}

Se $x_k$ è l'approssimazione disponibile all'iterazione $k$,
definiamo il residuo come

\begin{equation*}
    r_k=b-Ax_k.
    \label{eq:residuo-implementazione}
\end{equation*}

Se $x_k$ fosse la soluzione esatta, avremmo

\begin{equation*}
    Ax_k=b
\end{equation*}

e pertanto

\begin{equation*}
    r_k=0.
\end{equation*}

Se invece $x_k$ è soltanto un'approssimazione, il residuo sarà in
generale diverso dal vettore nullo.

Possiamo quindi utilizzare la norma del residuo,

\begin{equation*}
    \|r_k\|,
\end{equation*}

come misura di quanto l'approssimazione corrente non soddisfa
ancora il sistema.

Più precisamente, possiamo considerare

\begin{equation*}
    \|r_k\|_2
    =
    \|b-Ax_k\|_2.
    \label{eq:norma-residuo}
\end{equation*}

Se il metodo converge, ci aspettiamo che questa quantità diventi
progressivamente più piccola.


\subsection{Tolleranza e criterio di arresto}

In un calcolo numerico non è generalmente necessario richiedere
che il residuo sia esattamente uguale a zero.

Definiamo invece una quantità positiva molto piccola,
che indichiamo con $\varepsilon$ e chiamiamo \emph{tolleranza}.

Possiamo considerare l'approssimazione sufficientemente accurata
quando

\begin{equation*}
    \|r_k\|_2
    \leq
    \varepsilon.
    \label{eq:stopping-rule}
\end{equation*}

Il procedimento deve quindi continuare mentre

\begin{equation*}
    \|r_k\|_2
    >
    \varepsilon.
\end{equation*}

Questa condizione si traduce naturalmente in un ciclo
\texttt{while}:

\begin{verbatim}
while (norm(r, type = "2") > tolleranza) {

    # nuova iterazione

}
\end{verbatim}

La tolleranza stabilisce quindi il livello di precisione richiesto.
Una tolleranza più piccola impone un criterio di arresto più
stringente e può richiedere un numero maggiore di iterazioni.


\subsection{Costruzione dell'algoritmo iterativo}

Possiamo ora combinare gli strumenti introdotti in questa sezione
con il meccanismo iterativo studiato nella sezione precedente.

Partiamo dal sistema

\begin{equation*}
    Ax=b.
\end{equation*}

Assegniamo un'approssimazione iniziale

\begin{equation*}
    x_0
\end{equation*}

e calcoliamo il residuo iniziale

\begin{equation*}
    r_0=b-Ax_0.
\end{equation*}

Finché la norma del residuo è superiore alla tolleranza,
continuiamo a calcolare nuove approssimazioni.

Lo schema generale dell'algoritmo può essere rappresentato come

\begin{verbatim}
scegli x0
calcola r = b - A x0

while (norma(r) > tolleranza) {

    calcola la correzione z
    aggiorna x
    calcola il nuovo residuo r

}
\end{verbatim}

In termini matematici, utilizzando il meccanismo introdotto
precedentemente, a ogni iterazione calcoliamo

\begin{equation*}
    r_k=b-Ax_k,
\end{equation*}

risolviamo

\begin{equation*}
    Mz_k=r_k
\end{equation*}

e aggiorniamo la soluzione secondo

\begin{equation*}
    x_{k+1}=x_k+z_k.
\end{equation*}

Il procedimento continua fino a quando

\begin{equation*}
    \|b-Ax_k\|_2
    \leq
    \varepsilon.
\end{equation*}


\subsection{Il rischio di un ciclo infinito}

L'utilizzo di un ciclo \texttt{while} richiede una particolare
attenzione.

Se il metodo iterativo non converge, oppure se la condizione di
arresto non può essere soddisfatta, il ciclo potrebbe continuare
indefinitamente.

Per evitare questo problema è opportuno introdurre anche un numero
massimo di iterazioni, che indichiamo con $K_{\max}$.

Il procedimento viene quindi interrotto quando si verifica almeno
una delle seguenti condizioni:

\begin{equation*}
    \|r_k\|_2\leq\varepsilon
\end{equation*}

oppure

\begin{equation*}
    k\geq K_{\max}.
\end{equation*}

La struttura generale diventa quindi

\begin{verbatim}
while (norm(r, type = "2") > tolleranza &&
       iterazione < max_iterazioni) {

    # calcolo della nuova approssimazione

}
\end{verbatim}

Il numero massimo di iterazioni costituisce una misura di
sicurezza: impedisce al programma di continuare indefinitamente
nel caso in cui il metodo non converga.


\subsection{Interpretazione dell'algoritmo}

Possiamo infine rappresentare l'intera procedura come

\begin{equation*}
\boxed{
\begin{array}{c}
\text{scegliere }x_0\\
\downarrow\\
r_0=b-Ax_0\\
\downarrow\\
\text{calcolare }\|r_0\|\\
\downarrow\\
\|r_k\|\leq\varepsilon\;?
\end{array}}
\end{equation*}

Se la risposta è positiva, il procedimento termina e $x_k$ viene
accettato come soluzione numerica.

Se la risposta è negativa, viene calcolata una nuova
approssimazione:

\begin{equation*}
\boxed{
x_k
\longrightarrow
r_k
\longrightarrow
z_k
\longrightarrow
x_{k+1}
\longrightarrow
r_{k+1}
}
\end{equation*}

e il criterio di arresto viene verificato nuovamente.

Questa struttura evidenzia il ruolo dei diversi strumenti
computazionali introdotti nella sezione:

\begin{itemize}
    \item il ciclo \texttt{for} permette di ripetere una procedura
          per un numero prefissato di volte;
    \item il ciclo \texttt{while} permette di ripeterla finché
          una determinata condizione rimane verificata;
    \item la norma permette di associare a un vettore una misura
          scalare della sua grandezza;
    \item il residuo misura quanto l'approssimazione corrente non
          soddisfa ancora il sistema;
    \item la tolleranza stabilisce quando l'approssimazione può
          essere considerata sufficientemente accurata.
\end{itemize}

Gli strumenti introdotti in questa sezione permettono quindi di
passare dalla formulazione matematica di un metodo iterativo alla
sua effettiva implementazione in un linguaggio di programmazione.

